% Irodov user PDF, printed pp.52-54, PDF pp.61-63,
% hydrodynamics problems 325-334, Figs 1.87-1.89 inspected.
\exproblem{325}{52}
从理想流体的欧拉运动方程出发，
证明在不可压缩流体的定常流动中，
沿任一流线成立伯努利方程。
\exsolution
理想流体在重力场中的欧拉方程为
\[
\rho\frac{D\bm v}{Dt}=-\nabla p+\rho\bm g.
\]
定常流动时
$\partial\bm v/\partial t=0$，因此
$D\bm v/Dt=(\bm v\cdot\nabla)\bm v$。
令流线弧长为 $s$、单位切向矢量
$\bm t=\bm v/v$，
把上式与 $\bm t$ 作点积：
\[
\rho v\frac{\mathrm dv}{\mathrm ds}
=-\frac{\mathrm dp}{\mathrm ds}
-\rho g\frac{\mathrm dz}{\mathrm ds}.
\]
这里竖直坐标 $z$ 向上，故
$\bm g\cdot\bm t=-g\,\mathrm dz/\mathrm ds$。
在 $\rho$ 为常数的不可压缩近似下，
\[
\frac{\mathrm d}{\mathrm ds}
\left(\frac12\rho v^2+p+\rho gz\right)=0.
\]
对同一条流线积分：
\[
\boxed{\frac12\rho v^2+p+\rho gz=C_{\rm streamline}}.
\]
\excheck 每条流线的常数一般可能不同；
只有特定条件（如无旋势流）才能进一步说明不同
流线拥有同一个伯努利常数。

\exproblem{326}{52}
一个很宽的竖直水槽两相对侧壁分别有两个
截面积相同的小孔 $S=0.50\,\mathrm{cm^2}$。
两孔高度差 $\Delta h=51\,\mathrm{cm}$。
打开后两股水流向相反水平方向喷出。
求容器受到的合反冲力大小。
\exsolution
设两孔位于水面下的深度分别为 $h_1,h_2$。
由托里拆利定律，两孔出口速率满足
$v_i^2=2gh_i$。
射流质量通量为
$\dot m_i=\rho S v_i$，
容器对应的反冲力模为
$F_i=\dot m_i v_i=\rho S v_i^2
=2\rho g S h_i$。
两股喷流方向相反，反冲合力大小为
\[
F_{\rm net}=|F_1-F_2|
=2\rho gS|h_1-h_2|
=\boxed{2\rho gS\Delta h}.
\]
换算 $S=5.0\times10^{-5}\,\mathrm{m^2}$、
$\Delta h=0.51\,\mathrm m$，
\[
\boxed{F_{\rm net}
=2(1000)(9.81)(5.0\times10^{-5})(0.51)
\simeq0.50\,\mathrm N}.
\]
\excheck 水槽越深的孔喷速越大、反冲力越大，
净反冲方向指向较深孔\emph{喷射方向的反方向}。

\exproblem{327}{52}
一个宽的竖直水槽高 $H=75\,\mathrm{cm}$，
侧壁上有一条自槽底向上延伸的狭窄竖直缝，
长度为 $\ell=50\,\mathrm{cm}$、宽
$b=1.0\,\mathrm{mm}$。
开缝前水槽装满水，求刚打开缝时
全部出流水对槽壁的合反冲力大小。
\exsolution
设自槽底向上量坐标 $y\in[0,\ell]$。
此处水面下深度为 $H-y$，
出口速率由托里拆利公式为
$v(y)=\sqrt{2g(H-y)}$。
微小出流面积 $\mathrm dA=b\,\mathrm dy$。
单位时间流出的水平动量为
$\mathrm dF=\rho v(y)^2\,\mathrm dA
=2\rho g(H-y)b\,\mathrm dy$。
沿整个缝长度积分：
\[
F=\int_0^\ell2\rho gb(H-y)\,\mathrm dy
=2\rho gb(H\ell-\ell^2/2).
\]
因此
\[
\boxed{F=\rho gb\ell(2H-\ell)}.
\]
代入数值：
$\boxed{F\simeq4.91\,\mathrm N\approx5.0\,\mathrm N}$。
\excheck 这是\emph{打开瞬间}的反冲力。
随水面下降，各点深度改变，反冲力不再保持不变。

\exproblem{328}{53}
一个大水箱经内径半径
$r=0.50\,\mathrm{cm}$ 的直角弯管出水
（原书图1.87）。
弯管水平部分长 $\ell=22\,\mathrm{cm}$，
相对点 $O$ 的水平动量作用力臂为 $\ell$。
体积流量为 $Q=0.50\,\mathrm{L/s}$。
求出流水对管壁关于 $O$ 的反作用\emph{力矩大小}。
\exsolution
设管截面积 $A=\pi r^2$。
由不可压缩连续性，出水速度
$v=Q/A$。垂直向下的出口质量流率
$\dot m=\rho Q$，每秒携出的竖向动量为
$\dot m\,v=\rho Q^2/A$。
相对入口点 $O$，垂直动量通量的力臂为
$\ell$；入口水平动量通量沿通过 $O$ 的方向，
不贡献该点的力矩。
因此液体对弯管的反力矩模为
\[
\boxed{N=\ell\frac{\rho Q^2}{\pi r^2}}.
\]
取水密度 $\rho=1000\,\mathrm{kg/m^3}$、
$Q=5.0\times10^{-4}\,\mathrm{m^3/s}$、
$r=5.0\times10^{-3}\,\mathrm m$、$\ell=0.22\,\mathrm m$，
\[
\boxed{N\simeq0.70\,\mathrm{N\,m}}.
\]
\excheck 必须将升/秒换算成立方米/秒；
该式假定忽略管内流动产生的复杂附加壁面力矩与
非均匀出口速度剖面。

\exproblem{329}{53}
一个大开口水箱侧壁装有向外收缩的水平短管
（原书图1.88），入口横截面积
$S=3.0\,\mathrm{cm^2}$、出口面积
$s=1.0\,\mathrm{cm^2}$。
自由水面高出短管中心 $h=4.6\,\mathrm m$，
忽略黏性。求试图将短管从水箱侧壁
\emph{向外拔出}的水平水动力大小。
\exsolution
出水口通大气，流速
$v_s=\sqrt{2gh}$，
入口流速由连续性为
$v_S=(s/S)v_s$。
入口相对大气的静压由伯努利式为
\[
p_S-p_0=\rho gh-\frac12\rho v_S^2
=\rho gh\left(1-\frac{s^2}{S^2}\right).
\]
取收缩短管内水为控制体。
其水平动量流出减流入的通量为
$\rho s v_s(v_s-v_S)$。
入口压强、出口大气压与短管壁对流体作用力
合起来提供该动量变化。
以相对大气压计，短管\emph{对水}的水平合力为
\[
F_{{\rm wall}\to{\rm water}}
=\rho s v_s(v_s-v_S)-(p_S-p_0)S
=-\rho gh\,\frac{(S-s)^2}{S}.
\]
牛顿第三定律使水对短管向外的作用力为
\[
\boxed{F_{\rm pull}=\rho gh\frac{(S-s)^2}{S}}.
\]
代入 SI 单位 $S=3\times10^{-4}\,\mathrm{m^2}$、
$s=10^{-4}\,\mathrm{m^2}$，得
$\boxed{F_{\rm pull}\simeq6.0\,\mathrm N}$。
\excheck 不可只用静水压强乘出口面积，
还必须计入收缩管两端的动量流率变化。

\exproblem{330}{53}
一个盛水的竖直圆柱形容器以恒角速度
$\omega$ 绕自身对称轴旋转，液体最终达到
相对容器静止的稳定状态。
求：(a) 自由液面的形状；
(b) 如果槽底中心静压为 $p_0$，
槽底离轴距离 $r$ 处的静压强。
\exsolution
在随容器转动的参考系中，
液体处于相对静力平衡：
\[
\nabla p=\rho(\omega^2r\,\bm e_r-g\bm e_z).
\]
因此径向与竖直偏导数满足
$\partial_rp=\rho\omega^2r$、
$\partial_zp=-\rho g$。
积分得
\[
p(r,z)=C+\tfrac12\rho\omega^2r^2-\rho gz.
\]
自由液面压强为大气常数，
故自由表面相对其轴心的高度为
\[
\boxed{z(r)-z(0)=\frac{\omega^2r^2}{2g}},
\]
是旋转抛物面。
槽底各点同高 $z=0$，
压强分布为
\[
\boxed{p_{\rm bottom}(r)=p_0+\frac12\rho\omega^2r^2}.
\]
\excheck 自由表面上压强处处相同，
但高度随 $r$ 增加；底面同高而压强随 $r$ 增加，
二者不矛盾。

\exproblem{331}{53}
半径为 $R=10\,\mathrm{cm}$ 的水平薄圆盘
置于充满油的水平圆柱形狭缝内
（原书图1.89）。
圆盘与上、下固定平面的间隙均为
$h=1.0\,\mathrm{mm}$，
油的动力黏度
$\eta=0.08\,\mathrm{P}$
（泊，$1\,\mathrm P=0.1\,\mathrm{Pa\,s}$）。
圆盘绕中心轴以角速度
$\omega=60\,\mathrm{rad/s}$ 旋转。
忽略边缘效应，求黏滞力耗散功率。
\exsolution
盘面距轴半径为 $r$ 处线速度为
$v_\theta=\omega r$。在窄隙内，近似线性速度剖面
产生剪切应力
\[
\tau(r)=\eta\frac{\omega r}{h}.
\]
单侧盘面半径 $r$、宽 $\mathrm dr$ 的环带面积
$2\pi r\,\mathrm dr$，
其摩擦力矩微元
$\mathrm dN_1=\tau(r)\,2\pi r\,\mathrm dr\cdot r
=2\pi\eta\omega r^3\,\mathrm dr/h$。
上下两面都有黏滞作用，总力矩
\[
N=2\int_0^R\mathrm dN_1
=\frac{\pi\eta\omega R^4}{h}.
\]
所需驱动功率等于力矩乘角速度：
\[
\boxed{P=N\omega=\frac{\pi\eta\omega^2R^4}{h}}.
\]
用 $\eta=0.008\,\mathrm{Pa\,s}$ 计算，
$\boxed{P\simeq9.0\,\mathrm W}$。
\excheck 必须包含圆盘的\emph{上下两个面}，
并正确把泊换算为 SI 黏度单位。

\exproblem{332}{53}
半径为 $R_1$ 的无限长内圆柱
以恒定轴向速度 $v_0$ 在
半径为 $R_2>R_1$ 的静止同轴外圆柱内移动，
中间充满黏性液体。
忽略端部效应、假定稳定层流且无轴向压强梯度。
求液体轴向速度 $v_z(r)$ 的径向分布。
\exsolution
流动只沿轴向，速度只随半径 $r$ 变化。
牛顿黏性流体的稳定轴向运动方程在无压强梯度时为
\[
0=\eta\frac1r\frac{\mathrm d}{\mathrm dr}
\left(r\frac{\mathrm dv_z}{\mathrm dr}\right).
\]
积分一次给 $r\,\mathrm dv_z/\mathrm dr=C_1$，
再积分得到
$v_z(r)=C_1\ln r+C_2$。
利用无滑移边界
$v_z(R_1)=v_0$、$v_z(R_2)=0$，
\[
\boxed{v_z(r)=v_0
\frac{\ln(r/R_2)}{\ln(R_1/R_2)}},
\qquad R_1\le r\le R_2.
\]
\excheck 在内壁为 $v_0$、外壁为零，
中间连续单调下降；圆柱坐标中的分布是对数函数，
不是两平行平板间的严格线性函数。

\exproblem{333}{53--54}
两长同轴圆柱之间有黏度为 $\eta$ 的层流液体，
内柱半径 $R_1$ 且静止，
外柱半径 $R_2>R_1$、以角速度 $\omega_2$
绕公共轴旋转。已知半径 $r$ 处圆柱面上的剪切应力
$\tau(r)=\eta r\,\mathrm d\omega_f/\mathrm dr$，
其中 $\omega_f(r)$ 是流体局部角速度。
求：(a) $\omega_f(r)$；
(b) 单位轴向长度上作用于外柱的摩擦力矩大小。
\exsolution
稳定轴对称层流中，单位长度通过半径 $r$ 圆柱面的
扭矩为
\[
N'(r)=2\pi r^2\tau(r)
=2\pi\eta r^3\frac{\mathrm d\omega_f}{\mathrm dr}.
\]
稳态角动量传递要求 $N'(r)$ 与 $r$ 无关。
故 $\mathrm d\omega_f/\mathrm dr=C/r^3$，
积分得到
$\omega_f(r)=A+B/r^2$。
无滑移边界为
$\omega_f(R_1)=0$、
$\omega_f(R_2)=\omega_2$。
解得
\[
\boxed{\omega_f(r)=
\frac{\omega_2R_2^2}{R_2^2-R_1^2}
\left(1-\frac{R_1^2}{r^2}\right)}.
\]
由导数
$\omega_f'(r)=
2\omega_2R_1^2R_2^2/
[(R_2^2-R_1^2)r^3]$，
得到外柱受到的阻转力矩每单位长度大小
\[
\boxed{N'=
\frac{4\pi\eta\omega_2 R_1^2R_2^2}
{R_2^2-R_1^2}}.
\]
\excheck 该结果是\emph{单位轴向长度}的力矩，
量纲为 $\mathrm{N\,m/m}=\mathrm N$；
实际长为 $L$ 的外柱总阻转矩为 $N'L$。

\exproblem{334}{54}
半径为 $R$、长度为 $\ell$ 的圆管中流过
密度为 $\rho$、动力黏度为 $\eta$ 的稳定层流，
轴向流速剖面为
$v(r)=v_0(1-r^2/R^2)$。
求：(a) 体积流量；
(b) 管内流体动能；
(c) 流体对管壁的总黏滞摩擦力；
(d) 两端压强差。
\exsolution
(a) 以同心环带积分流量：
\[
Q=\int_0^R v(r)2\pi r\,\mathrm dr
=2\pi v_0\int_0^R(r-r^3/R^2)\,\mathrm dr
=\boxed{\frac12\pi v_0R^2}.
\]
(b) 管内各环带长度均为 $\ell$，动能
\[
K=\frac12\rho\ell\int_0^R v(r)^2\,2\pi r\,\mathrm dr
=\boxed{\frac16\pi\ell R^2\rho v_0^2}.
\]
(c) 牛顿黏性应力在管壁的大小
$|\tau_w|=\eta|v'(R)|=2\eta v_0/R$。
壁面积 $2\pi R\ell$，所以
\[
\boxed{F_f=|\tau_w|(2\pi R\ell)=4\pi\eta\ell v_0}.
\]
(d) 稳态充分发展流动中，
端面压强推动力平衡壁面黏滞阻力，
$\Delta p\,(\pi R^2)=F_f$，故
\[
\boxed{\Delta p=\frac{4\eta\ell v_0}{R^2}}.
\]
\excheck 与泊肃叶定律
$Q=\pi R^4\Delta p/(8\eta\ell)$
互相一致；本题 $v_0$ 是轴心最大速率，
平均速率为 $v_0/2$。
